\section{Generación de carga desde los clientes}
\label{sec:carga}

\subsection{Herramienta}

Se usará un \textbf{generador de carga propio en C++}: el mismo programa cliente
(\codigo{bin/cliente}) con un modo de carga que reutiliza la clase \codigo{Cliente}.
Se descartaron herramientas como JMeter, Locust o k6 porque hablan HTTP y no
ZeroMQ; usarlas obligaría a escribir un adaptador que agregaría su propio tiempo a
las mediciones. Con el cliente propio, la carga recorre exactamente el mismo camino
que un estudiante real.

\begin{lstlisting}
bin/cliente --carga consulta|consulta-reserva|reserva-modificacion
            --duracion 720 --id C1 --escenario ESC-5 --arquitectura base
            --csv resultados/tiempos_ESC-5_base_C1.csv
\end{lstlisting}

\subsection{Método}

\begin{itemize}
    \item \textbf{Un proceso por cliente simulado}, cada uno con su propio socket
          \codigo{REQ}. Así los clientes son realmente concurrentes y no comparten
          colas internas.
    \item \textbf{Ciclo cerrado sin tiempo de espera:} cada cliente envía la siguiente
          solicitud apenas recibe la respuesta. Con pocos clientes (3 a 8) esto mide
          la capacidad real del sistema y da resultados repetibles.
    \item \textbf{Tipos de cliente} según la Tabla~1:
          \begin{itemize}
              \item \emph{consulta}: repite \codigo{consultarEventosMes} sobre los meses con eventos;
              \item \emph{consulta-reserva}: una consulta y luego una reserva de 1 cupo;
              \item \emph{reserva-modificacion}: una reserva y luego su modificación
                    (cambio de cupos).
          \end{itemize}
    \item \textbf{Datos de prueba:} cada reserva la hace un estudiante distinto,
          generado para la prueba, en un evento con capacidad suficiente. Así no se
          choca con la regla de una reserva por evento ni se agotan los cupos, y los
          errores de negocio no contaminan los tiempos.
    \item \textbf{Orquestación:} un script \codigo{bash},
          \codigo{scripts/correr\_escenario.sh}, recibe el escenario
          (\codigo{ESC-3}, \codigo{ESC-5} o \codigo{ESC-8}) y la arquitectura
          (\codigo{base} o \codigo{escalable}), lanza los procesos del escenario al mismo tiempo, espera a que terminen y
          une sus archivos CSV.
    \item \textbf{Ubicación:} los clientes corren en una máquina distinta a la de los
          servicios (Máquina 1 de la Figura~\ref{fig:despliegue}) para no competir por
          CPU con el sistema medido.
    \item \textbf{Pruebas de estrés y resistencia:} el mismo generador, lanzando más
          procesos de forma escalonada.
\end{itemize}

Para la prueba de concurrencia C-01 se lanza un grupo de procesos que envían su
reserva al mismo tiempo sobre una fecha con pocos cupos.
